% Generated draft: requires mathematical and source review.
\documentclass[11pt]{article}
\usepackage[margin=25mm]{geometry}
\usepackage{fontspec}
\setmainfont{Latin Modern Roman}
\usepackage{amsmath,amssymb,mathtools,unicode-math}
\setmathfont{Latin Modern Math}
\usepackage{enumitem,needspace,xurl,hyperref}
\setlength{\parindent}{0pt}
\setlength{\parskip}{5pt}
\setlength{\emergencystretch}{3em}
\newcommand{\N}{\mathbb{N}}
\newcommand{\Z}{\mathbb{Z}}
\newcommand{\Q}{\mathbb{Q}}
\newcommand{\R}{\mathbb{R}}
\newcommand{\C}{\mathbb{C}}
\newcommand{\F}{\mathbb{F}}
\newcommand{\A}{\mathbb{A}}
\newcommand{\PP}{\mathbb{P}}
\newcommand{\E}{\mathbb{E}}
\newcommand{\eps}{\varepsilon}
\newcommand{\dd}{\,\mathrm{d}}
\begin{document}
\section*{UnsolvedMath ID: 30004637}
\textbf{Source problem code:} catalogue entry
\subsection{Definitions and mathematical statement}
The supplied record asks about a regularized unbalanced optimal-transport problem associated with branching Brownian motion. Here the branching law, Brownian state space and diffusivity, time horizon, admissible transport data, optimization variables, objective functional, regularization, and boundary conditions are unspecified. Consequently, the record does not determine a unique mathematical optimization problem.

The supplied algorithmic question is: \emph{Can one design a fast algorithm for the regularized unbalanced optimal-transport problem associated with branching Brownian motion?}

The meaning of ``fast'' is also unspecified: the record gives no input representation, computational model, accuracy criterion, convergence requirement, or complexity bound. A precise quantified algorithmic statement therefore requires additional specifications from the underlying source. No particular formulation or performance requirement is imposed in this draft.
\subsection{Short English statement}
The question asks for a fast algorithm for regularized unbalanced transport associated with branching Brownian motion. The supplied record does not specify the optimization problem or the required computational performance, so the intended algorithmic target remains unresolved.
\subsection{Sources}
\begin{itemize}
\item Local problem registry: \texttt{data/unsolvedmath/problems.json}, permanent UnsolvedMath ID \(30004637\), source problem code \texttt{catalogue entry}, in the supplied collection \emph{Oberwolfach Reports Open Problems}. No individual proposer is identified; the supplied proposed year is 2021.
\item Supplied report citation: \emph{Applications of Optimal Transportation in the Natural Sciences} (2022), Oberwolfach Reports, pp.~507--587, DOI \texttt{10.4171/owr/2021/10}. Supplied source and literature link: \url{https://doi.org/10.4171/owr/2021/10}, labeled in the literature sources as ``Oberwolfach Report 10/2021.'' The supplied background attributes the problem to mathematical workshop open-problem discussions documented by Oberwolfach Reports.
\item Supplied literature assessment, checked on exactly \texttt{2026-08-22}: the record labels the problem open and states that the cited report gives no algorithmic resolution and that the general problem remains open. This is supplied dated curation, described in the background as literature triage requiring independent verification of time-sensitive status and cited claims; it has not been independently verified for this draft.
\end{itemize}
\end{document}
